%; whizzy paragraph -pdf xpdf -latex ./whizzypdfptex.sh
%; whizzy-paragraph "^\\\\begin{frame}\\|\\\\emtext"
% latex beamer presentation.
% platex, latex-beamer でコンパイルすることを想定。 

%     Tokyo Debian Meeting resources
%     Copyright (C) 2012 Junichi Uekawa
%     Copyright (C) 2011 Nobuhiro Iwamatsu

%     This program is free software; you can redistribute it and/or modify
%     it under the terms of the GNU General Public License as published by
%     the Free Software Foundation; either version 2 of the License, or
%     (at your option) any later version.

%     This program is distributed in the hope that it will be useful,
%     but WITHOUT ANY WARRANTY; without even the implied warreanty of
%     MERCHANTABILITY or FITNESS FOR A PARTICULAR PURPOSE.  See the
%     GNU General Public License for more details.

%     You should have received a copy of the GNU General Public License
%     along with this program; if not, write to the Free Software
%     Foundation, Inc., 51 Franklin St, Fifth Floor, Boston, MA  02110-1301 USA

\documentclass[cjk,dvipdfmx,12pt]{beamer}
\usetheme{Tokyo}
\usepackage{monthlypresentation}

%  preview (shell-command (concat "evince " (replace-regexp-in-string "tex$" "pdf"(buffer-file-name)) "&")) 
%  presentation (shell-command (concat "xpdf -fullscreen " (replace-regexp-in-string "tex$" "pdf"(buffer-file-name)) "&"))
%  presentation (shell-command (concat "evince " (replace-regexp-in-string "tex$" "pdf"(buffer-file-name)) "&"))

%http://www.naney.org/diki/dk/hyperref.html
%日本語EUC系環境の時
\AtBeginDvi{\special{pdf:tounicode EUC-UCS2}}
%シフトJIS系環境の時
%\AtBeginDvi{\special{pdf:tounicode 90ms-RKSJ-UCS2}}

\title{東京エリアDebian勉強会}
\subtitle{第87回 2012年4月度}
\author{野島 貴英\\nozzy@debian.or.jp}
\date{2012年4月21日}
\logo{\includegraphics[width=8cm]{image200607/openlogo-light.eps}}

\begin{document}

\frame{\titlepage{}}

\begin{frame}{設営準備にご協力ください。}
会場設営、会誌ホチキス止めなどよろしくおねがいします。
\end{frame}

\section{お願い}
\begin{frame}{お願い}
\begin{itemize}
\item お願い
  \begin{itemize}
  \item セミナールーム、トイレ、喫煙室以外は立ち入らないように願います。
  \item 写真撮影はセミナールームの中のみでお願いします。
  \end{itemize}
\end{itemize}
\end{frame}

\section{Agenda}
\begin{frame}{Agenda}
  \begin{itemize}
  \item 最近あったDebian関連のイベント報告
    \begin{itemize}
    \item 第85回 東京エリアDebian勉強会
    \item 第86回 東京エリアDebian勉強会
    \end{itemize}
  \end{itemize}
\end{frame}

\section{Agenda(つづき)}
\begin{frame}{Agenda(つづき)}
\begin{itemize}
\item DWN quiz
\item 事前課題紹介
\item Debianでのnode入門
\item Android機でDebian
\item 月刊Debhelper
\item 今後のイベント
\item 今日の宴会所
\end{itemize}
\end{frame}
\begin{frame}{第85回 東京エリアDebian勉強会}
\begin{itemize}
\item 開催場所は三栄町生涯学習館 視聴覚室(JR四ツ谷駅近く)
\item KDEパッケージ開発 / Debhelper(dh\_dpatch\_patch,dh\_autotools-dev\_updateconfig) / cmakeの話題を行いました。
\item 参加者5人でした。プログラマ向けの話題でもり上がりました。
\end{itemize}
\end{frame}

\begin{frame}{第86回 東京エリアDebian勉強会}
\begin{itemize}
\item 開催場所は明星大学
\item テーマは「Debian 勉強会 -ユーザサイド- Apache2 HTTP サーバから始めるDebian」というユーザサイドの発表が行われました。
\item Debian JP Projectとしてブースも出しました。Debian Live CDも配りました。あんどきゅめんてぃっど でびあんもいくつか売れていました。
\end{itemize}
\end{frame}

\section{DWN quiz}
\emtext{DWN quiz}
\begin{frame}{Debian 常識クイズ}

Debian の常識、もちろん知ってますよね?
知らないなんて恥ずかしくて、知らないとは言えないあんなことやこんなこと、
みんなで確認してみましょう。

今回の出題範囲は\url{debian-devel-announce@lists.debian.org},
\url{debian-devel@lists.debian.org} に投稿された
内容とDebian Project Newsなどからです。

\end{frame}

\subsection{問題}
%; whizzy-master ../debianmeetingresume201204.tex
% $B0J>e$N@_Dj$r$7$F$$$k$?$a!"$3$N%U%!%$%k$G(B M-x whizzytex $B$9$k$H!"(Bwhizzytex$B$,MxMQ$G$-$^$9!#(B
%

\santaku
{$B:#G/EY$N(BDebian JP$B2qD9$OC/$+!)(B}
{Kouhei Maeda}
{Nobuhiro Iwamatsu}
{Junichi Uekawa}
{A}
{$BA0G/EY$K0z$-B3$-A0ED$5$s$,?.G$$5$l$^$7$?!#$^$?$h$m$7$/$*4j$$$7$^$9!#(B}


\santaku
{Debian.org DPL $BA*5s$OC/$,N)8uJd$7$?$+(B}
{Stefano Zacchiroli}
{Nobuhiro Iwamatsu}
{Kouhei Maeda}
{A}
{$BA0G/EY$K0z$-B3$-(BZacchiroli$B$5$s4hD%$C$F$^$9!#(B}

\santaku
{Debconf12$B$N(Bsuponsord$B$J;22C$NDy@Z$j$O$$$D(B}
{4$B7nKv(B}
{5$B7n(B15$BF|(B}
{5$B7nKv(B}
{B}
{$B$&$C$+$j$9$k$H2a$.$F$7$^$&$N$G!"5$$r$D$1$^$7$g$&!#$^$?!"(BUTC$B$J$N$+%K%+%i%0%"$N%m!<%+%k%?%$%`$J$N$+$A$g$C$H$o$+$i$J$$$N$G!"(B5/15$B$KEPO?3+;O$9$k$N$OHr$1$?J}$,$h$$$+$b!#K\2H%"%J%&%s%9(B\url{http://debconf12.debconf.org/}$B;2>H!#(B}

\santaku
{$BBgE}0l(BDebian$BJY6/2q$G$NH/I=$N8xJg(B(CFP)$B$O$$$D$,Dy@Z$j!)(B}
{$B$b$&2a$.$?(B}
{4$B7nKv(B}
{4$B7n(B22$BF|(B}
{C}
{$BBgE}0l(BDebian$BJY6/2q$GH/I=$G$-$k%A%c%s%9$G$9!#Dy@Z$j$OK:$l$:$K!#(B}

\santaku
{experimental$BHG$N(Bapache$B$N%Q%C%1!<%8$N%P!<%8%g%s$O$$$/$D!)(B}
{2.3}
{2.4.0}
{2.4.2}
{C}
{3/22$B$K%"%J%&%s%9$,$"$j$^$7$?!#(Bupstream$BB&$b(B2.4.2$B$G$9!#:G?7HG$G$9$M!#(BBUG$B8+$D$1$^$7$?$i!"(BBUG$B!!(BReport$B=q$-$^$7$g$&!#(B}

\santaku
{Debian Edu$B$O$^$?$NL>$r2?$H$$$&$G$7$g$&(B?}
{emdebian}
{Scientific Linux}
{Skolelinux}
{C}
{Debian Edu$B$H$O650i5!4X8~$1$K:n$i$l$?(BDebian$B%Y!<%9$N%G%#%9%H%j%S%e!<%7%g%s$N;v$G$9!#@N$O(BSkolelinux$B$H$$$&L>A0$G3+H/$5$l$F$$$?J*$@$=$&$G$9!#@5<0$J(BDebian$B$N%5%V%W%m%8%'%/%H$G$9!#@hF|?7$7$$%P!<%8%g%s$,%"%J%&%s%9$5$l$^$7$?!#(B}

\santaku
{3/30$B$K(BDebian Project$B$,2CLA$7$?CDBN$NL>A0$O!)(B}
{OSC}
{OSI}
{ETF}
{B}
{$B$3$l$G$^$?(BFree Software$B$H$7$FKa$-$,$+$+$j$^$7$?!#(BOSI$B$O(BOpen Source Initiative$B$NN,$@$=$&$G$9!#(B}

\santaku
{Debian Project$B$N(Bgobby$B%5!<%P!<$H$7$F(Bgobby.debian.org$B$,%"%J%&%s%9$5$l$^$7$?!#$H$3$m$G(Bgobby$B$C$F2?!)(B}
{$B5l%=O"$G3+H/$5$l$?D5Js3hF0MQ%=%U%H%&%'%"(B}
{churro$B$NBeBX%5!<%P(B}
{$B%(%G%#%?$NL>A0(B}
{C}
{Debconf$B$N(BBOF$B2q>l$G$h$/;H$o$l$F$$$^$9%(%G%#%?$G$9!#%5!<%P$r2p$9$k;v$K$h$j!"J#?t?M$GF1;~$K#1$D$NJ8>O$rF1;~$KJT=8$G$-$^$9!#(BDebconf$B$G$O!"%j%"%k%?%$%`$K5D;v$H5D;vO?$,(BBOF$B;22C<T$K$h$C$F$I$s$I$sJT=8$5$l$F$$$/MM$O$*$b$7$m$$$G$9!#$A$J$_$K(Bchurro$B$H$O(Bi18n.debian.org$B$N;v$G$9!#(B}




\section{事前課題}
\emtext{事前課題}
{\footnotesize
 \begin{prework}{ koedoyoshida }
\begin{enumerate}
\item  Debian$B$G?($C$?$3$H$N$"$k%&%'%V%"%W%j%1!<%7%g%s%U%l!<%`%o!<%/0l$D$r4J7i$K>R2p$7$F$/$@$5$$!#(B\\
Greasemonkey
\item $B6a!9:n@.$7$F$_$?$$$H;W$C$F$$$k%&%'%V%"%W%j%1!<%7%g%s$r65$($F$/$@$5$$!#(B\\
$B$J$7(B 
\end{enumerate}
\end{prework}

\begin{prework}{ noritadak }
\begin{enumerate}
\item Ruby on Rails$B!#$G$b$+$J$j@N$G$9!#(B
\item Web $B%"%W%j%1!<%7%g%s$H8@$($k$N$+HyL/$G$9$,!"%V%i%&%6>e$G;H$($k%7%'%k$r8D?ME*$K:n$j$?$$$G$9!#(B
\end{enumerate}
\end{prework}

\begin{prework}{ dictoss($B?yK\!!E5=<(B) }
\begin{enumerate}
\item python$B$N(BDjango$B$O$J$+$J$+$$$$46$8$G$9$,(Bpython3$B$K8=>uHsBP1~$H$$$&$H$3$m!#(Bcakephp$B$b$[$s$N>/$7;n$7$^$7$?$,$h$/$o$i$+$::C@^!#(B
\item $B<+Bp$NK\$r4IM}$9$k%7%9%F%`$,$[$7$$!#$^$?!"$[$7$$%8%c%s%k$NK\$r<+F0$GC5$7$F$/$l$k%\%C%H$D$-$,$$$$$J!#(B
\end{enumerate}
\end{prework}

\begin{prework}{ henrich }
\begin{enumerate}
\item $B%&%'%V%"%W%j%1!<%7%g%s%U%l!<%`%o!<%/!"0l$D$b?($C$?$3$H$,$4$6$$$^$;$s!D(B
\item $B:n@.$7$F$_$?$$$H$$$&$N$OFC$K$OL5$$$G$9$,!">-MhE*$K$O(BDjango$BJU$j$r?($C$F$_$?$$$H$O;W$C$F$$$^$9!#(B
\end{enumerate}
\end{prework}

\begin{prework}{ ikeomasa }
\begin{itemize}
\item django
\end{itemize}
\end{prework}

\begin{prework}{ MATOHARA }
\begin{enumerate}
\item Debian$B$G?($C$?$3$H$N$"$k%&%'%V%"%W%j%1!<%7%g%s%U%l!<%`%o!<%/0l$D$r4J7i$K>R2p$7$F$/$@$5$$!#(B \\
$B%U%l!<%`%o!<%/$OMxMQ$7$?$3$H$,$"$j$^$;$s!%(BPerl $B$N(BCatalyst $B$r;H$C$F$_$?$$$J$H;W$C$F$$$^$9!%(B
\item $B6a!9:n@.$7$F$_$?$$$H;W$C$F$$$k%&%'%V%"%W%j%1!<%7%g%s$r65$($F$/$@$5$$!#(B\\
$B:G6a<L??$N(Bexif $B>pJs$rJT=8$9$k$b$N$r:n$j$?$$$J$H;W$C$F$$$^$9!%(B\\
$B!t%"%C%W%m!<%I$7$?<L??$N(Bexif $B>pJs$rI=<($7$F%+%a%i$N%7%j%"%k!$;#1F<T!$0LCV>pJs$J$I$rA*Br$7$F:o=|$7$F%@%&%s%m!<%I!%(B
\end{enumerate}
\end{prework}

\begin{prework}{ beatenavenue }
\begin{enumerate}
\item pukiwiki$B$b(Bweb$B%"%W%j$KF~$j$^$9$G$7$g$&$+!&!&!&!#%0%k!<%W%&%'%"$_$?$$$J;H$$J}$,$G$-$J$$$+$J$H;W$C$F;n$7$F$$$^$7$?$,!"Cg4V$&$A$G(Bwiki$BJ8K!$NI>H=$,Hs>o$K0-$/:#$O?($C$F$$$^$;$s!#(B
\item $B:n@.$7$?$$(Bweb$B%"%W%j$O:#$N$H$3$m$"$j$^$;$s!#(B
\end{enumerate}
\end{prework}

\begin{prework}{ $BK\>1(B }
\begin{enumerate}
\item CakePHP$B$r;H$C$?$3$H$,$"$j$^$9!#(Bpear$B$K0MB8$7$J$$$H$$$&$3$H$GA*Br$7$^$7$?!#(B
\item $B:G6a!"%F%l%SHVAH$r%A%'%C%/$9$k<+J,MQ$NE,Ev$J%9%/%j%W%H$r:n@.$7$^$7$?!#(B
\end{enumerate}
\end{prework}

\begin{prework}{ yamamoto }
\begin{enumerate}
\item $BA4$/?($C$?$3$H$,$"$j$^$;$s!#(B
\item $B%&%'%V%"%W%j%1!<%7%g%s$C$F!"$J$s$9$+$M!)$*$$$7$$$N!)(B
\end{enumerate}
\end{prework}

\begin{prework}{ $B$^$($@$3$&$X$$(B }
\begin{enumerate}
\item Sinatra $B;H$C$F$^$9!#(BDebian$B$8$c$J$/$F(BUbuntu $B$G$9$,!D!#(B
\item $B6a!9!"$H$$$&$+(B1$B$G(BOS$B%$%s%9%H!<%kMQ%D!<%k$N3+H/$d$C$F$^$9!#(B
\end{enumerate}
\end{prework}

\begin{prework}{ $BNkLZ?rJ8(B }
\begin{enumerate}
\item $B=q@R$K$N$C$F$$$?%5%s%W%k$rF0:n$5$;$k$?$a$K(Bpylons$B$r>/$7$@$1;H$C$?$3$H$,$"$j$^$9!#8=:_$O(Bpyramid$B$rBe$o$j$K;HMQ$9$k$3$H$,?d>)$5$l$F$$$k$_$?$$$G$9!#$H$3$m$G!"%U%l!<%`%o!<%/$r;H$C$F(Bweb$B3+H/$9$k>l9g$O$I$&$d$C$F(Bdebian$B%Q%C%1!<%8$rMxMQ$9$Y$-$J$s$G$7$g$&$+!#$$$^$^$G$O%Q%C%1!<%8$r;HMQ$;$:$KD>@\%@%&%s%m!<%I$7$F;H$C$F$$$^$7$?!#(B
\item $B0LCV>pJs$r;H$C$?(BWeb$B%"%W%j$r:n$C$F$_$?$$$G$9!#(B
\end{enumerate}
\end{prework}

\begin{prework}{ emasaka }
\begin{enumerate}
\item Bash on Rails$B:n$C$?$o!<(B 3$BG/$0$i$$A0$K:n$C$?$o!<!J%_%5%o!K(B\\
$B!J$H2]Bj$K=q$/$?$a$K;22C$7$^$9!K(B
\item $B;E;v$G$A$g$C$H$7$?(BWeb$B%"%W%j$rDs0FCf(B
\end{enumerate}
\end{prework}

\begin{prework}{ $BBg?9!!=S=((B }
\begin{enumerate}
\item django
\item $B$+$s$?$s=5Js%"%W%j(B
\end{enumerate}
\end{prework}
\begin{prework}{ $BLnEg!!5.1Q(B }
\begin{enumerate}
\item $B8E$$OC$G$9$,!"(BPerl$B$N(BCatalyst$B$0$i$$$G$7$g$&$+(B?($B:#$@$H(Bmojiolicious$B$J$s$@$m$&$+(B...$B!#(B)$B>R2p$K$D$$$F$O(B...$B$Y!"JY6/$7$F$-$^$9!<!#(B
\item $B%5%$%H=d2s$7$F2?$+%G!<%?<h$C$F$-$F<+F0E*$KJ,N`!?;w$F$$$k$b$N$rC5$98D?MMxMQ@lMQ$N%5%$%H$r:n$C$F$_$?$$$J$!(B...$B$"!"$3$l$@$H%&%'%V%U%l!<%`%o!<%/$h$j$b!"5!3#3X=,$H$+$N%F!<%^(B...$B$4$U$C(B...
\end{enumerate}
\end{prework}

}


\section{Debian での node 入門}
\emtext{Debian での node 入門}

\begin{frame}{Debian でNode を使う}
\begin{itemize}
 \item nodejs
 \item nodejs-dev
 \item npm
\end{itemize}
\end{frame}

\begin{frame}{Debian node policy}
 \begin{itemize}
  \item node-dev README.Debian
  \item モジュールのパッケージ名はnode-で始まる
  \item モジュールのロードパスの規程
 \end{itemize}
\end{frame}

\begin{frame}{npm / Debian package path}
\begin{table}[h]
\begin{center}
 \begin{tabular}{|l|l|}
 \hline
 & ディレクトリ\\
 \hline
 Debian パッケージ & \texttt{/usr/lib/nodejs} \\
 Debian の npm  & \texttt{/usr/local/lib/nodejs}\\
 npm install -g & \texttt{/usr/lib/node\_{}modules} \\
 npm install & \texttt{./node\_{}modules} \\
 \hline
 \end{tabular} 
\end{center}
\end{table}
\end{frame}

\begin{frame}{npm}
\debianbug{622628} is now fixed.
\end{frame}

\begin{frame}[containsverbatim]{npm}
\begin{commandline}
$ npm help
$ npm help npm
$ npm help install
\end{commandline}
\end{frame}

\begin{frame}[containsverbatim]{npm}
 \begin{commandline}
$ npm install -g express ejs
$ npm link
 \end{commandline}
\end{frame}

\begin{frame}[containsverbatim]{package.json}
 man npm-json.1\\
\begin{commandline}
{
    "name": "aptserver",
    "version": "v0.0.1",
    "dependencies": {
	"cli": "",
	"express": "",
	"ejs": ""
    }
} 
\end{commandline}
\end{frame}

%$
\section{Android機でDebian}
\emtext{Android機でDebian}


\begin{frame}{モバイル機/タブレットでDebian}

\center{モバイル機/タブレットでDebian
が動いたらいいとおもいませんか？}

\end{frame}

\begin{frame}{ところでみなさんスマホもってますか？}

\center{ところでみなさんスマホもってますかーっ？}

\end{frame}

\begin{frame}{そのうちAndroid使ってますかー}

\center{うちAndroid使ってますかー？}

\end{frame}

\begin{frame}{そのうちroot奪取は？}

\center{root奪取してますかー？}

\end{frame}

\begin{frame}{AndroidSDK入れてadb使ってますかー}

\center{AndroidSDK入れてadb使ってますかー}

※俺の環境じゃー動かんけどなー？

\end{frame}

\begin{frame}{簡単にDebian（派生も？）とか入れるには}
今だとroot奪取してアプリでできるらしい（すみません自分は未評価です）
\begin{itemize}
\item ubuntuを動作させるAndroid アプリ：ubuntu instller free \url{https://play.google.com/store/apps/details?id=com.zpwebsites.ubuntuinstall&hl=ja}
\item Debianを動作させるAndroid アプリ: Lil' Debi \url{https://github.com/guardianproject/lildebi/wiki}
\end{itemize}
\end{frame}

\begin{frame}{今回の材料：Nook Color}
\begin{itemize}
\item Barnes \& Noble社の電子書籍端末
\item 軽い、薄い
\item LCD、
\item タッチパネル
\item Wi-Fi
\item miniSD I/F
\item USBポート
\item Android OS搭載（かなりカスタマイズされています）
\end{itemize}
\end{frame}

\begin{frame}{Nook Colorのもっと凄いところ}
\begin{itemize}
\item TI社OMAP3チップ(ARM CPUですよ)搭載
\item miniSDからブートできちゃう。
\end{itemize}
\end{frame}

\begin{frame}{Debianの凄いところ}
\begin{itemize}
\item ARMでもたくさんのパッケージが動いちゃう。
\item muti-archと言われるパッケージはx86上でもARM用のツールにつかえちゃう。\\
例: objdumpとか！x86の元で逆ASMしたいですよね！？当然ですよね！？
\item emdebianのパッケージ入れると、クロスコンパイル環境も作れちゃう！
\item ARMのエミュレータだってパッケージでさっくり入る
\end{itemize}
\end{frame}

\begin{frame}{さあ、「インストールだ！」}

さあ、「インストールだ！」\\
 qemu-system-armで。\\
※詳しくは月刊Deb専2012年4月21日を参照
\end{frame}

\begin{frame}{さあ、「動け！！」}

さあ、「動け！！」\\
 qemu-system-armで。\\
※詳しくは月刊Deb専2012年4月21日を参照
\end{frame}

\begin{frame}{目覚めよ!OMAP3}
目覚めよ!OMAP3ということでブートフォーマットについて記載します。
\begin{itemize}
\item microSDのHEAD 255,SECTOR 63,512[BYTE/SEC]を使ってパーティション切りをし、
\item 第一パーティションをvfatで32bitのファットサイズにしておく
\item あとはTI提供のx-bootをMLOとかにリネームして、カーネルイメージと、initrdと、u-bootファイルとMLOファイルをそのままコピーしておくだけ。
\end{itemize}
以上を満たすmicroSDがあればブートしちゃいます。
\end{frame}

\begin{frame}{カーネルと起動イメージの入手}
nooter0.2.zipから借用します。\\
理由：usb経由のRNDIS超便利ーっ\\
※詳しくは月刊Deb専2012年4月21日を参照
\end{frame}

\begin{frame}{デモ}
デモ\\
※作成済みのminiSDから起動して、GIMP/Youtube HTML5動画再生まで。
\end{frame}


\section{月刊Debhelper}
\emtext{月刊Debhelper}

\section{今後のイベント}
\emtext{今後のイベント}
\begin{frame}{今後のイベント}
\begin{itemize}
\item 大統一Debian勉強会\\
日時：6月23日(土曜日)\\
場所：京都大学 理学研究科3号館 数学教室 108, 109, 110 号室
\end{itemize}
\end{frame}

\section{今日の宴会場所}
\emtext{今日の宴会場所}
\begin{frame}{今日の宴会場所}
\end{frame}

\end{document}

;;; Local Variables: ***
;;; outline-regexp: "\\([ 	]*\\\\\\(documentstyle\\|documentclass\\|emtext\\|section\\|begin{frame}\\)\\*?[ 	]*[[{]\\|[]+\\)" ***
;;; End: ***
